\documentclass[11pt,a4paper]{article}

\usepackage[utf8]{inputenc}
\usepackage[margin=1in]{geometry}
\usepackage{amsmath,amssymb}
\usepackage{booktabs}
\usepackage{graphicx}
\usepackage{microtype}
\usepackage{cite}
\usepackage{xcolor}
\usepackage[colorlinks=true,linkcolor=blue,citecolor=blue,urlcolor=blue]{hyperref}

% Every number below comes from tables/numbers.tex, generated by make_tables.py
% from results/v2 (harness v2). Do not type numbers into the text by hand.
\newcommand{\benefitAllZeroPct}{58}
\newcommand{\cmpCycCNNCIvsCycCNN}{-0.29}
\newcommand{\cmpCycCNNCIvsCycCNNwins}{17/28}
\newcommand{\cmpCycCNNCIvsCycLinCI}{+15.13}
\newcommand{\cmpCycCNNCIvsCycLinCIwins}{1/32}
\newcommand{\cmpCycFixedvsLinear}{-2.27}
\newcommand{\cmpCycFixedvsLinearwins}{27/32}
\newcommand{\cmpCycLinCIvsCycLin}{-1.94}
\newcommand{\cmpCycLinCIvsCycLinwins}{17/28}
\newcommand{\cmpCycLinCIvsLinear}{-1.75}
\newcommand{\cmpCycLinCIvsLinearwins}{26/32}
\newcommand{\cmpCycMLPCIvsCycLinCI}{-1.94}
\newcommand{\cmpCycMLPCIvsCycLinCIwins}{20/32}
\newcommand{\cmpCycMLPCIvsCycMLP}{-1.64}
\newcommand{\cmpCycMLPCIvsCycMLPwins}{25/28}
\newcommand{\cmpCycSORCIvsCycLinCI}{+0.52}
\newcommand{\cmpCycSORCIvsCycLinCIwins}{14/32}
\newcommand{\cmpCycSORCIvsCycSOR}{-1.99}
\newcommand{\cmpCycSORCIvsCycSORwins}{20/28}
\newcommand{\cmpDCTLinvsLinear}{-0.33}
\newcommand{\cmpDCTLinvsLinearwins}{24/32}
\newcommand{\cmpDLinearvsLinear}{-0.12}
\newcommand{\cmpDLinearvsLinearwins}{23/32}
\newcommand{\cmpHaarLinvsLinear}{-0.16}
\newcommand{\cmpHaarLinvsLinearwins}{25/32}
\newcommand{\cmpPatchTSTvsCycFixed}{-0.55}
\newcommand{\cmpPatchTSTvsCycFixedwins}{20/32}
\newcommand{\cmpPatchTSTvsCycMLPCI}{-0.92}
\newcommand{\cmpPatchTSTvsCycMLPCIwins}{20/32}
\newcommand{\cmpSTFTLinvsLinear}{+1.95}
\newcommand{\cmpSTFTLinvsLinearwins}{4/32}
\newcommand{\cmprFFTLinvsLinear}{-0.30}
\newcommand{\cmprFFTLinvsLinearwins}{25/32}
\newcommand{\cnnParamsETT}{5{,}559}
\newcommand{\consensusAllZeroPct}{83}
\newcommand{\cycDailyGainSolar}{-3.5}
\newcommand{\cycLinEtthTwoSevenTwenty}{0.4298}
\newcommand{\dailyFVUEhOne}{0.42}
\newcommand{\dailyFVUEhTwo}{0.61}
\newcommand{\dailyFVUEmTwo}{0.58}
\newcommand{\dailyFVUSolar}{0.24}
\newcommand{\dailyFVUTraffic}{0.18}
\newcommand{\dailyFVUWeather}{0.75}
\newcommand{\dailyShareEhOne}{50}
\newcommand{\dailyShareEhTwo}{12}
\newcommand{\dailyShareEmTwo}{18}
\newcommand{\dailyShareSolar}{57}
\newcommand{\dailyShareTraffic}{46}
\newcommand{\dailyShareWeather}{23}
\newcommand{\dlinearEtthTwoSevenTwenty}{0.3931}
\newcommand{\dsofCNNbenefitMed}{+0.0}
\newcommand{\dsofCNNbenefitWins}{3/8}
\newcommand{\dsofLINEARhardMed}{-0.0}
\newcommand{\dsofLINEARhardWins}{4/8}
\newcommand{\dsofSORbenefitMed}{+0.0}
\newcommand{\dsofSORbenefitWins}{3/8}
\newcommand{\dsofSORconsensusMed}{+0.0}
\newcommand{\dsofSORconsensusWins}{1/8}
\newcommand{\dsofSORhardMed}{-5.0}
\newcommand{\dsofSORhardWins}{8/8}
\newcommand{\dsofSORsoftMed}{-3.4}
\newcommand{\dsofSORsoftWins}{7/8}
\newcommand{\elecCycFixed}{0.1582}
\newcommand{\elecCycLinCI}{0.1690}
\newcommand{\fdecChronosBoltImproved}{1/32}
\newcommand{\fdecChronosTFiveImproved}{13/32}
\newcommand{\fdecTTMImproved}{2/8}
\newcommand{\fdecTimesFMTwoFiveImproved}{0/32}
\newcommand{\fmChronosBoltVsLinMed}{+12.7}
\newcommand{\fmChronosBoltVsLinWins}{3/32}
\newcommand{\fmChronosBoltVsSnWins}{28/32}
\newcommand{\fmChronosTFiveVsLinMed}{+81.9}
\newcommand{\fmChronosTFiveVsLinWins}{0/32}
\newcommand{\fmChronosTFiveVsSnWins}{14/32}
\newcommand{\fmTimesFMTwoFiveVsLinMed}{+7.7}
\newcommand{\fmTimesFMTwoFiveVsLinWins}{7/32}
\newcommand{\fmTimesFMTwoFiveVsSnWins}{31/32}
\newcommand{\fourierBaseWorseThanSnaive}{18/32}
\newcommand{\levelShareEhOne}{16}
\newcommand{\levelShareEhTwo}{39}
\newcommand{\levelShareEmTwo}{48}
\newcommand{\levelShareSolar}{24}
\newcommand{\levelShareTraffic}{7}
\newcommand{\levelShareWeather}{46}
\newcommand{\levslowShareEhOne}{39}
\newcommand{\levslowShareEhTwo}{78}
\newcommand{\levslowShareEmTwo}{74}
\newcommand{\levslowShareSolar}{37}
\newcommand{\levslowShareTraffic}{22}
\newcommand{\levslowShareWeather}{69}
\newcommand{\linFamSpreadBigData}{1.1}
\newcommand{\linFamSpreadMed}{1.0}
\newcommand{\linParamsETT}{32{,}534}
\newcommand{\medGapFifth}{2.7}
\newcommand{\medGapSecond}{0.6}
\newcommand{\medSeedCV}{0.44}
\newcommand{\medWithinThree}{5.5}
\newcommand{\mlpLevelGainSolar}{-6.7}
\newcommand{\mlpLevelGainWeather}{-7.6}
\newcommand{\mlpParamsETT}{221{,}974}
\newcommand{\nCells}{32}
\newcommand{\nDistinctWinners}{6}
\newcommand{\nModels}{19}
\newcommand{\nRuns}{1696}
\newcommand{\nSeedCells}{544}
\newcommand{\nTrained}{17}
\newcommand{\nUnstable}{18}
\newcommand{\parsevalMaxErr}{4e-16}
\newcommand{\patchDailyGainSolar}{-11.6}
\newcommand{\patchLevelGainSolar}{-7.4}
\newcommand{\patchLevelGainWeather}{-8.0}
\newcommand{\rankCycCNNCI}{15.9}
\newcommand{\rankCycFixed}{5.2}
\newcommand{\rankCycLinCI}{5.9}
\newcommand{\rankCycMLPCI}{5.3}
\newcommand{\rankDLinear}{8.9}
\newcommand{\rankLinear}{9.8}
\newcommand{\rankPatchTST}{3.9}
\newcommand{\rankSTFTLin}{13.0}
\newcommand{\rankrFFTLin}{7.1}
\newcommand{\rfftEtthTwoSevenTwenty}{0.3911}
\newcommand{\sorControlVsWarmMed}{+2.2}
\newcommand{\sorHardVsLinCtrlMed}{+3.3}
\newcommand{\timesfmTrafficWins}{4}
\newcommand{\trafficSevenTwentyCycLin}{0.6986}
\newcommand{\trafficSevenTwentyCycLinStd}{0.2013}
\newcommand{\ttmVsLinWins}{1/8}
\newcommand{\winsCycFixed}{8}
\newcommand{\winsCycMLPCI}{6}
\newcommand{\winsPatchTST}{11}
\newcommand{\winsTF}{5}
\newcommand{\xmAgree}{311/312}
\newcommand{\xmMedAbs}{0.000}


\newcommand{\lrn}[1]{\textbf{L#1}}
\newcommand{\todo}[1]{\textcolor{red}{[TODO: #1]}}

\title{\textbf{When Does Frequency-Domain Modeling Help in Lightweight Forecasting?\\Lessons from a Controlled Study of Inductive Biases}}

\author{
    \textbf{Aaditya Chaturvedi}\\
    Information Systems and Machine Learning Lab (ISMLL)\\
    University of Hildesheim, Germany\\
    \texttt{chaturvedia@uni-hildesheim.de}
}
\date{\today}

\begin{document}
\maketitle

\begin{abstract}
Lightweight multivariate forecasters increasingly add frequency-domain components: Fourier or wavelet bases, learned periodic states, spectral gating, Fourier-based decoupling of pretrained models, and spectral distillation. Each is usually reported with a new architecture and a favorable table, which makes it hard to tell which ingredient does the work. We instead hold one protocol fixed and vary the ingredients. On eight standard datasets, four horizons ($L{=}336$, $H\in\{96,192,336,720\}$), three seeds and per-cell learning-rate selection, we train \nTrained{} models (\nRuns{} runs, reproduced across two machines with \xmAgree{} cells agreeing within seed spread) and evaluate four zero-shot foundation models. We decompose every error exactly into level, slow ($>$36\,h), daily (8--36\,h) and fast ($<$8\,h) bands. The main lessons are: (i) among strong lightweight models the ranking is narrow and dataset dependent (median gap between the best and fifth-best model \medGapFifth\%, six different winners over 32 cells); (ii) a fixed orthonormal transform (rFFT, DCT, DWT) in front of a linear head changes accuracy by only $0.1$--$0.3\%$, because it does not change what the model can represent; (iii) an explicit periodic term helps when the period is real and stable (median $\cmpCycLinCIvsLinear\%$ vs.\ a plain linear map, \cmpCycLinCIvsLinearwins{} cells) and hurts on the drifting ETTh2, while choosing the wrong period costs more than the term gains; (iv) on drifting datasets most error sits in the level and slow bands, where no model in the study beats a constant (FVU $\approx 1$); (v) the largest gains (PatchTST, MLP residual on Solar and Weather) come mostly from the level band, not from periodic structure; (vi) zero-shot foundation models beat seasonal naive but rarely a trained linear model, Fourier decoupling of them mostly hurts, and spectral distillation helps only as a regularizer of an overfitting student. We describe the corrected diagnostics and the claims of our earlier draft that did not survive.
\end{abstract}

\section{Introduction}
\label{sec:intro}
Since DLinear~\cite{zeng2023dlinear} showed that a linear map can match Transformer forecasters on the standard long-horizon benchmarks~\cite{wu2021autoformer,zhou2022fedformer,nie2023patchtst}, a large part of the field has moved to small models with hand-designed inductive biases. Many of them are spectral: frequency-domain projections~\cite{zhou2022fedformer}, wavelet mixing~\cite{wpmixer2024}, explicit learned cycles~\cite{liu2024cyclenet,lin2025tqnet}, and frequency-aware losses~\cite{FreqCycle2026,FreqLens2026}. Usually each idea comes as part of one new architecture, tuned by its authors and compared with numbers copied from earlier papers under different settings. When the claimed margins are a few percent, protocol differences of the same size make it unclear what works.

This project started the same way: we built a model (CycleNet-Pro, an extension of CycleNet~\cite{liu2024cyclenet} with automatic period discovery and four residual backbones) and a spectral diagnostic, and reported favorable results. A later audit found that several of those results came from protocol differences, single seeds and diagnostics that do not mean what we said (Section~\ref{sec:removed}). We therefore re-ran everything under one pre-registered protocol and turned the question around. Instead of asking whether \emph{our} model wins, we ask:

\begin{quote}
\emph{When does frequency-domain modeling actually help, and when does it fail, in lightweight multivariate forecasting?}
\end{quote}

We answer it by isolating ingredients: fixed spectral transforms, explicit periodic states, nonlinear and patch-based residual networks, channel mixing, zero-shot foundation models with and without Fourier decoupling, and spectral distillation. All of them use the same data pipeline, training recipe, metric code and seeds. Our contributions are lessons, not a new model:
\begin{itemize}
    \item A controlled benchmark of \nTrained{} trained models and four zero-shot models on 32 dataset--horizon cells, with three seeds, learning-rate selection per cell, and a cross-machine reproduction check (Section~\ref{sec:protocol}).
    \item A corrected spectral diagnostic: an exact Parseval split of MSE into \emph{physically defined} bands, plus skill against seasonal naive (Section~\ref{sec:diagnostics}).
    \item Eleven findings (\lrn{1}--\lrn{11}, Section~\ref{sec:results}), each with its evidence and its caveat, including several negative results.
    \item An explicit list of earlier claims that we withdraw and why (Section~\ref{sec:removed}).
\end{itemize}

\section{Related Work}
\label{sec:related}
\paragraph{Linear and lightweight forecasters.} DLinear~\cite{zeng2023dlinear} decomposes the lookback into a moving-average trend and a remainder and maps each linearly to the horizon. RevIN~\cite{kim2021revin} normalizes each input window and restores the statistics at the output; nearly all recent lightweight models use it. PatchTST~\cite{nie2023patchtst} uses channel-independent patching with a Transformer encoder; HDMixer~\cite{hdmixer2024} replaces attention with hierarchical patch mixing.

\paragraph{Explicit periodicity.} CycleNet~\cite{liu2024cyclenet} learns a cycle table of length $p$ per channel, subtracts the phase-aligned cycle from the input and adds it back to the forecast, letting a small backbone model only the residual. TQNet~\cite{lin2025tqnet} extends this with temporal queries for cross-variable dependencies.

\paragraph{Spectral representations.} FEDformer~\cite{zhou2022fedformer} works with a subset of Fourier and wavelet modes; WPMixer~\cite{wpmixer2024} mixes wavelet packet coefficients; WaveMoE~\cite{WaveMoE2026} routes tokens across wavelet subbands. Recent work argues that MSE training concentrates fit in low frequencies~\cite{FreqCycle2026,FreqLens2026}.

\paragraph{Foundation models.} TimesFM~\cite{das2024timesfm} is a decoder-only patch model; Chronos~\cite{ansari2024chronos} tokenizes values for a T5 language model, and Chronos-Bolt maps patches directly to the horizon; TTM~\cite{ekambaram2024ttm} is a small pretrained MLP-mixer.

\paragraph{Evaluation.} We report skill against the seasonal naive forecast, in the spirit of MASE~\cite{hyndman2006another}. For claimed pairwise winners we use paired tests across cells; Diebold--Mariano tests~\cite{diebold1995comparing} on per-window losses are left for the next revision (Section~\ref{sec:threats}).

\section{Protocol and Model Families}
\label{sec:protocol}

\subsection{Protocol}
All numbers come from one harness (\texttt{srp26/bench}, v2), with these rules fixed before the runs:
\begin{itemize}
    \item \textbf{Data.} ETTh1, ETTh2, ETTm1, ETTm2 (7 channels; 12/4/4-month chronological split), Electricity (321), Solar (137), Traffic (862), Weather (21) (70/10/20\% split). The scaler is fitted on the training split only; validation and test windows include the preceding $L$ steps; dataset checksums are verified. Metrics are computed in scaled space.
    \item \textbf{Task.} Lookback $L{=}336$; horizons $H\in\{96,192,336,720\}$; stride-1 test windows.
    \item \textbf{Training.} Adam, cosine schedule, FP32, MSE loss, batch 32, at most 30 epochs, early stopping on validation MSE (patience 5). The learning rate is chosen per (dataset, horizon, model) on validation MSE with seed 2024 from $\{10^{-4},5\cdot10^{-4},10^{-3},5\cdot10^{-3}\}$; if the best value is on the edge of the grid, the grid is extended ($5\cdot10^{-5}$, $10^{-5}$ or $10^{-2}$). PatchTST uses its published recipe (fixed $10^{-4}$, up to 100 epochs, patience 10).
    \item \textbf{Seeds and reproduction.} Final runs use seeds 2024, 2025 and 2026 (mean $\pm$ std reported). The ETT part of the grid was run a second time on a different machine and GPU; \xmAgree{} model--cell pairs agree within seed spread (median absolute relative difference \xmMedAbs\%).
\end{itemize}

\subsection{Model families}
Each family tests one inductive bias; all trained models except PatchTST share the RevIN wrapper.

\paragraph{References (untrained).} Seasonal naive (repeat the last day; 24, 96 or 144 steps), repeat-last.

\paragraph{Linear maps with fixed transforms.} \emph{Linear} maps the raw lookback to the horizon per channel. \emph{DLinear} applies the moving-average decomposition first. The \emph{TF-} models apply a fixed transform to the lookback and then a linear head: rFFT (real and imaginary parts), orthonormal DCT-II, two-level orthonormal DWT (Haar, db4), and an overlapping Hann-window STFT. Because every transform here is linear in the input, all of these models represent the same set of functions as Linear; only the parameterization, and with it the optimization, differs. This family therefore isolates the effect of the \emph{representation} from the effect of the \emph{function class}.

\paragraph{Explicit periodic state (CycleNet family).} A learned cycle table of length $p$ per channel is subtracted from the input and added to the forecast, and a backbone models the residual. \emph{Cyc-Fixed} uses the cycle length from the CycleNet paper (24, 96, 168 or 144 steps) with a linear backbone. \emph{CycleNet-Pro} chooses the periods automatically (FFT peaks of the channel-mean training series, filtered by autocorrelation significance and a variance-reduction rule), decides automatically whether to add a channel-mixing MLP, and offers four backbones: linear, two-layer MLP, a small 1D-CNN, and a spectrally conditioned patch mixer (SOR). The \texttt{-CI} variants force channel independence. We treat CycleNet-Pro as one family among others.

\paragraph{Patch Transformer.} PatchTST/42, re-implemented without residual attention, with the official small configuration on ETTh1/2 and the large one elsewhere.

\paragraph{Zero-shot foundation models.} TimesFM-2.5 (200M), Chronos-Bolt-Base, Chronos-T5-Tiny and TTM-R1 (only $H{=}96$), applied to the same scaled lookbacks. Because inference is expensive, they are scored on non-overlapping test windows (stride $H$, covering the test period once) with one seed; trained models are re-scored on the same windows for comparison.

\paragraph{Fourier decoupling.} For each foundation model, the top-3 rFFT components of the lookback are extrapolated in closed form, the model forecasts the remainder, and the two are added. The top-3 base alone is also scored.

\paragraph{Spectral distillation (DSOF-SD).} A CycleNet-Pro student (linear, SOR or CNN backbone) is warm-started for 5 epochs and then trained for up to 25 epochs with an added loss that pulls its band spectrum towards the TimesFM forecast in selected bands (hard, soft, benefit-probe or consensus gates). Every gate is paired with a \emph{control}: the same trainer, schedule and data order with distillation weight zero. Distillation has to beat the control, not the warm student. This was run on ETTh1 and ETTh2.

\section{Diagnostics}
\label{sec:diagnostics}
\paragraph{Exact band decomposition.} For a forecast error $e\in\mathbb{R}^H$ with one-sided DFT $E_k$ and the usual weights ($w_0{=}1$, $w_{H/2}{=}1$ for even $H$, otherwise 2), Parseval gives $\frac1H\sum_t e_t^2=\frac1{H^2}\sum_k w_k|E_k|^2$. Assigning bins to disjoint bands gives an exact split
\begin{equation}
\mathrm{MSE}=\mathrm{MSE}_{\text{level}}+\mathrm{MSE}_{\text{slow}}+\mathrm{MSE}_{\text{daily}}+\mathrm{MSE}_{\text{fast}},
\end{equation}
verified numerically on every run (largest absolute deviation \parsevalMaxErr). The bands are defined by \emph{period in hours}, not by normalized frequency: level is the DC bin (error in the window mean), slow is periods above 36\,h, daily is 8--36\,h, fast is below 8\,h. With these bands the daily cycle and its first harmonics fall into the daily band at every sampling rate. When a band holds no bins at a given $H$ (for example slow at $H{=}96$ for 10- or 15-minute data) it is reported as missing.

\paragraph{Band FVU.} $\mathrm{FVU}_b$ divides the band error energy by the band energy of the target around its window mean. $\mathrm{FVU}_b<1$ means the model tracks some of the variation in that band. $\mathrm{FVU}_b\approx 1$ is what an MSE-optimal forecast produces in a band whose content is unpredictable from the lookback: shrinking towards zero there is correct behavior, not a failure. FVU therefore shows \emph{where predictability exists}, not which model is wrong.

\paragraph{Skill.} $\text{skill}=1-\mathrm{MSE}/\mathrm{MSE}_{\text{S-Naive}}$ against the seasonal naive forecast of the dataset's daily period.

\paragraph{What we no longer use.} We dropped three diagnostics from the earlier draft after testing them on synthetic cases: log-spectral distance as a quality score (an MSE-optimal smooth forecast scored 56.1\,dB and the same forecast with added fake noise scored 4.7\,dB while its MSE doubled, so lower LSD can mean a worse forecast); counts of ``failure windows'' (the threshold depended on the model's own error, and the count correlated 0.88 with MSE); and spectral coherence (with one Welch segment it equals 1 even against pure noise). These numbers come from the metric audit of 18 November 2026.

\section{Results: What We Learned}
\label{sec:results}

\begin{table*}[t]
\centering
\caption{Test MSE averaged over the four horizons (seed means), and average rank over the 32 cells among the 17 trained models. Bold: best trained model; underline: second best. Full per-horizon results in Table~\ref{tab:full}.}
\label{tab:avg}
\resizebox{\textwidth}{!}{\begin{tabular}{lccccccccccccc}
\toprule
Dataset & S-Naive & Linear & DLinear & rFFT-Lin & DCT-Lin & Haar-Lin & STFT-Lin & Cyc-Fixed & Cyc-Lin-CI & Cyc-MLP-CI & Cyc-SOR-CI & Cyc-CNN-CI & PatchTST \\
\midrule
ETTh1 & 0.600 & 0.417 & 0.424 & 0.418 & \textbf{0.414} & \underline{0.416} & 0.426 & 0.417 & 0.417 & 0.426 & 0.442 & 0.445 & 0.419 \\
ETTh2 & 0.483 & 0.347 & 0.347 & \underline{0.345} & \textbf{0.343} & 0.345 & 0.354 & 0.359 & 0.359 & 0.359 & 0.360 & 0.394 & 0.351 \\
ETTm1 & 0.489 & 0.365 & 0.363 & 0.360 & 0.360 & 0.367 & 0.372 & 0.357 & 0.357 & \underline{0.351} & 0.359 & 0.395 & \textbf{0.351} \\
ETTm2 & 0.358 & 0.263 & 0.263 & 0.257 & 0.256 & 0.257 & 0.268 & \textbf{0.254} & \underline{0.254} & 0.262 & 0.264 & 0.307 & 0.257 \\
Electricity & 0.330 & 0.169 & 0.169 & 0.169 & 0.169 & 0.169 & 0.176 & \textbf{0.158} & 0.169 & \underline{0.162} & 0.171 & 0.240 & 0.162 \\
Solar & 0.341 & 0.261 & 0.261 & 0.260 & 0.260 & 0.261 & 0.270 & 0.251 & 0.251 & 0.217 & \underline{0.211} & 0.278 & \textbf{0.200} \\
Traffic & 1.161 & 0.433 & 0.433 & 0.433 & 0.433 & 0.433 & 0.441 & 0.424 & \underline{0.422} & 0.422 & 0.428 & 0.607 & \textbf{0.421} \\
Weather & 0.371 & 0.249 & 0.247 & 0.248 & 0.248 & 0.249 & 0.259 & 0.244 & 0.244 & \underline{0.225} & 0.229 & 0.285 & \textbf{0.223} \\
\midrule
Avg.\ rank (of 17) & -- & 9.8 & 8.9 & 7.1 & 7.2 & 8.4 & 13.0 & 5.2 & 5.9 & 5.3 & 8.6 & 15.9 & 3.9 \\
\bottomrule
\end{tabular}}
\end{table*}

\paragraph{\lrn{1}: Among strong lightweight models the ranking is narrow and depends on the dataset.}
\emph{Evidence.} Over the 32 cells, the median gap between the best and the second-best trained model is \medGapSecond\%, and between the best and the fifth best \medGapFifth\%; a median of \medWithinThree{} models are within 3\% of the best. Six different models win at least one cell: PatchTST \winsPatchTST, Cyc-Fixed \winsCycFixed, Cyc-MLP-CI \winsCycMLPCI, and the transform-linear models \winsTF{} (Table~\ref{tab:avg}). PatchTST has the best average rank (\rankPatchTST), followed by Cyc-Fixed (\rankCycFixed) and Cyc-MLP-CI (\rankCycMLPCI). PatchTST wins mainly on Solar and Weather, the CycleNet variants on Electricity and Traffic, and the transform-linear models on ETTh2.
\emph{Caveat.} The median seed-to-seed coefficient of variation is \medSeedCV\%, so most gaps are larger than seed noise, but \nUnstable{} of \nSeedCells{} model--cell pairs have a coefficient of variation above 5\%. A single-seed leaderboard would reorder many cells.
\emph{Lesson.} A claim that one lightweight architecture is best ``overall'' is not supported at this scale of differences; claims should name the dataset regime.

\begin{table*}[t]
\centering
\caption{Paired comparisons over cells (seed-mean MSE). $\Delta=\mathrm{MSE}_A/\mathrm{MSE}_B-1$; negative favors A. Channel-mixing rows exclude ETTh1, where the automatic rule disabled mixing and both variants are identical. $p$-values are descriptive: cells of the same dataset are correlated.}
\label{tab:paired}
\resizebox{\textwidth}{!}{\begin{tabular}{lllcccc}
\toprule
Question & A & B & cells & median $\Delta$ (\%) & A better & Wilcoxon $p$ \\
\midrule
Fixed linear transform vs.\ raw lookback & rFFT-Lin & Linear & 32 & -0.30 & 25/32 & 0.001 \\
 & DCT-Lin & Linear & 32 & -0.33 & 24/32 & 0.0001 \\
 & Haar-Lin & Linear & 32 & -0.16 & 25/32 & 0.005 \\
 & DLinear & Linear & 32 & -0.12 & 23/32 & 0.04 \\
 & STFT-Lin & Linear & 32 & +1.95 & 4/32 & 3e-07 \\
Cycle queue (auto period) vs.\ none & Cyc-Lin-CI & Linear & 32 & -1.75 & 26/32 & 0.003 \\
Cycle queue (dataset period) vs.\ none & Cyc-Fixed & Linear & 32 & -2.27 & 27/32 & 0.0007 \\
Nonlinear residual (MLP) vs.\ linear & Cyc-MLP-CI & Cyc-Lin-CI & 32 & -1.94 & 20/32 & 0.05 \\
Patch mixer (SOR) vs.\ linear & Cyc-SOR-CI & Cyc-Lin-CI & 32 & +0.52 & 14/32 & 0.9 \\
1D-CNN vs.\ linear & Cyc-CNN-CI & Cyc-Lin-CI & 32 & +15.13 & 1/32 & 1e-09 \\
Channel independent vs.\ auto-mixing & Cyc-Lin-CI & Cyc-Lin & 28 & -1.94 & 17/28 & 0.4 \\
 & Cyc-MLP-CI & Cyc-MLP & 28 & -1.64 & 25/28 & 1e-06 \\
 & Cyc-SOR-CI & Cyc-SOR & 28 & -1.99 & 20/28 & 0.001 \\
 & Cyc-CNN-CI & Cyc-CNN & 28 & -0.29 & 17/28 & 0.7 \\
PatchTST vs.\ best CycleNet variant & PatchTST & Cyc-MLP-CI & 32 & -0.92 & 20/32 & 0.004 \\
 & PatchTST & Cyc-Fixed & 32 & -0.55 & 20/32 & 0.08 \\
\bottomrule
\end{tabular}}
\end{table*}

\paragraph{\lrn{2}: A fixed spectral transform in front of a linear head does not change what is learned.}
\emph{Evidence.} rFFT-, DCT- and Haar-linear models differ from the plain Linear model by a median of $\cmprFFTLinvsLinear\%$, $\cmpDCTLinvsLinear\%$ and $\cmpHaarLinvsLinear\%$ (Table~\ref{tab:paired}); DLinear by $\cmpDLinearvsLinear\%$. The overcomplete STFT is worse ($\cmpSTFTLinvsLinear\%$, better in only \cmpSTFTLinvsLinearwins{} cells). Within the linear family (Linear, DLinear, rFFT, DCT, both DWTs) the median spread per cell is \linFamSpreadMed\% and at most \linFamSpreadBigData\% on the three large datasets; the larger spreads on ETT come from unstable seeds (e.g.\ DLinear on ETTh1 $H{=}720$), not from the transform.
\emph{Explanation.} All these transforms are linear and (except STFT) orthonormal, so transform-then-linear is a reparameterized linear map. The small consistent advantage of rFFT/DCT is plausibly an optimization effect (better conditioning under a fixed epoch budget).
\emph{Lesson.} ``Working in the frequency domain'' is not by itself an inductive bias. A spectral component can only help if it changes the function class (e.g.\ a non-linear operation on coefficients) or the optimization target.

\begin{table}[t]
\centering
\caption{Per-dataset effects (relative MSE change in \%, averaged over horizons; negative is better). Cycle queue: Cyc-Lin-CI vs.\ Linear. Weekly vs auto period: Cyc-Fixed vs.\ Cyc-Lin-CI (identical where the periods agree). MLP residual: Cyc-MLP-CI vs.\ Cyc-Lin-CI. Channel mixing: Cyc-MLP (auto mixing) vs.\ Cyc-MLP-CI. PatchTST: vs.\ Cyc-Lin-CI. $C$: channels; $p$: periods used.}
\label{tab:effects}
\resizebox{\columnwidth}{!}{\begin{tabular}{lcllccccc}
\toprule
Dataset & $C$ & auto $p$ & fixed $p$ & Cycle queue & Weekly vs auto period & MLP residual & Channel mixing (auto) & PatchTST \\
\midrule
ETTh1 & 7 & 24 & 24 & -0.2 & +0.0 & +2.2 & +0.0 & +0.5 \\
ETTh2 & 7 & 24 & 24 & +3.2 & +0.0 & +0.5 & +0.9 & -2.0 \\
ETTm1 & 7 & 96 & 96 & -2.1 & +0.0 & -1.9 & +1.4 & -1.8 \\
ETTm2 & 7 & 16, 24, 96 & 96 & -3.6 & -0.0 & +3.8 & +1.2 & +1.8 \\
Electricity & 321 & 24 & 168 & -0.0 & -6.5 & -4.4 & +5.2 & -4.0 \\
Solar & 137 & 144 & 144 & -3.6 & -0.0 & -13.5 & +0.3 & -20.4 \\
Traffic & 862 & 168, 24 & 168 & -2.5 & +0.3 & -0.1 & +5.7 & -0.4 \\
Weather & 21 & 144 & 144 & -1.7 & +0.0 & -9.3 & +3.0 & -9.8 \\
\bottomrule
\end{tabular}}
\end{table}

\paragraph{\lrn{3}: An explicit periodic term helps when the period is real and stable, and the choice of period matters more than the term.}
\emph{Evidence.} Adding a learned cycle to a linear model (Cyc-Lin-CI vs.\ Linear) lowers MSE by a median of $\cmpCycLinCIvsLinear\%$ (\cmpCycLinCIvsLinearwins{} cells); with the dataset's known period (Cyc-Fixed) by $\cmpCycFixedvsLinear\%$ (\cmpCycFixedvsLinearwins). The gain is largest on Solar, ETTm2 and Traffic and appears in the daily band (Solar: $\cycDailyGainSolar$ points of MSE, Table~\ref{tab:bandgains}). On ETTh2 the cycle term \emph{hurts} by about 3\% (Table~\ref{tab:effects}). On Electricity the automatic rule picks the daily period (24) while the weekly period (168) is better: Cyc-Fixed reaches \elecCycFixed{} average MSE against \elecCycLinCI{}, a 6.5\% difference that is larger than the whole cycle effect on that dataset.
\emph{Caveat.} The automatic discovery runs on the channel-mean series and keeps peaks that pass an autocorrelation test, which favors the strongest short cycle; it found extra periods on ETTm2 (16, 24, 96) and Traffic (168, 24) without a clear benefit over the single fixed period.
\emph{Lesson.} The periodic prior is useful but cheap to get wrong. A period-selection rule should be validated on held-out error, not on spectral peak strength.

\begin{table}[t]
\centering
\caption{Where the error lives (Cyc-Lin-CI, averaged over horizons). Share of MSE per band, band FVU, and skill against seasonal naive. Missing bands at $H{=}96$ are skipped in the average.}
\label{tab:bands}
\resizebox{\columnwidth}{!}{\begin{tabular}{lccccccccc}
\toprule
& \multicolumn{4}{c}{Share of MSE (\%)} & \multicolumn{3}{c}{FVU} & Skill vs. \\
\cmidrule(lr){2-5}\cmidrule(lr){6-8}
Dataset & level & slow & daily & fast & slow & daily & fast & S-Naive (\%) \\
\midrule
ETTh1 & 16 & 23 & 50 & 10 & 1.05 & 0.42 & 0.85 & 30 \\
ETTh2 & 39 & 39 & 12 & 10 & 0.98 & 0.61 & 1.02 & 26 \\
ETTm1 & 23 & 19 & 50 & 13 & 0.95 & 0.36 & 0.78 & 27 \\
ETTm2 & 48 & 26 & 18 & 14 & 0.98 & 0.58 & 0.92 & 30 \\
Electricity & 26 & 28 & 30 & 15 & 0.63 & 0.07 & 0.49 & 49 \\
Solar & 24 & 13 & 57 & 12 & 0.76 & 0.24 & 0.60 & 26 \\
Traffic & 7 & 15 & 46 & 32 & 0.61 & 0.18 & 0.71 & 64 \\
Weather & 46 & 23 & 23 & 19 & 0.96 & 0.75 & 1.01 & 35 \\
\bottomrule
\end{tabular}}
\end{table}

\paragraph{\lrn{4}: Frequency modeling ``fails'' where the error is in the level and slow bands, and no model in this study predicts those.}
\emph{Evidence.} On ETTh2, ETTm2 and Weather, level and slow error together make up \levslowShareEhTwo\%, \levslowShareEmTwo\% and \levslowShareWeather\% of the MSE, and the slow-band FVU is close to 1 (Table~\ref{tab:bands}). On ETTh1, Solar and Traffic the daily band dominates instead (\dailyShareEhOne\%, \dailyShareSolar\%, \dailyShareTraffic\%) and is well predicted (daily FVU \dailyFVUEhOne, \dailyFVUSolar, \dailyFVUTraffic). The same slow-band FVU near 1 holds for PatchTST and the other models.
\emph{Interpretation.} The boundary is not a property of periodic models: on the drift-dominated datasets, the part of the future that carries most of the error is not predictable from a 336-step window by any model here. A periodic prior can only improve the daily band, which there holds only 12--23\% of the error.
\emph{Lesson.} Before adding a spectral component, measure how much of the baseline error lies in the band that component targets; that share bounds the possible gain.

\begin{table*}[t]
\centering
\caption{Where gains come from: change in band MSE as a percentage of the reference model's total MSE (negative is better; the four bands sum to the total change). Cycle queue: Cyc-Lin-CI minus Linear; MLP residual: Cyc-MLP-CI minus Cyc-Lin-CI; PatchTST: minus Cyc-Lin-CI.}
\label{tab:bandgains}
\resizebox{\textwidth}{!}{\begin{tabular}{lcccccccccccc}
\toprule
& \multicolumn{4}{c}{Cycle queue} & \multicolumn{4}{c}{MLP residual} & \multicolumn{4}{c}{PatchTST} \\
\cmidrule(lr){2-5}\cmidrule(lr){6-9}\cmidrule(lr){10-13}
Dataset & lev & slow & day & fast & lev & slow & day & fast & lev & slow & day & fast \\
\midrule
ETTh1 & +0.3 & +0.1 & -0.5 & -0.1 & +0.6 & -0.5 & +0.8 & +1.3 & -1.2 & -0.7 & +0.4 & +1.9 \\
ETTh2 & +2.0 & +0.6 & +0.1 & +0.5 & +0.1 & -0.5 & +0.9 & +0.1 & -2.1 & -0.2 & +0.3 & -0.0 \\
ETTm1 & -0.2 & -0.0 & -1.6 & -0.3 & -1.0 & +0.1 & -1.4 & +0.4 & -1.8 & -0.1 & -0.3 & +0.3 \\
ETTm2 & -1.1 & -0.8 & -1.3 & -0.6 & +1.7 & -0.0 & +1.7 & +0.4 & -0.8 & +0.3 & +1.4 & +1.1 \\
Electricity & +0.2 & +0.0 & -0.2 & -0.1 & -1.0 & -2.0 & -1.4 & +0.1 & -0.6 & -1.8 & -1.4 & -0.1 \\
Solar & -0.1 & -0.0 & -3.5 & +0.0 & -6.7 & +0.4 & -6.5 & -0.5 & -7.4 & +0.0 & -11.6 & -1.4 \\
Traffic & -0.0 & -0.9 & -1.1 & -0.3 & +0.6 & -0.0 & -0.3 & -0.3 & -0.5 & -0.4 & +0.6 & -0.2 \\
Weather & -1.5 & -0.2 & -0.7 & +0.7 & -7.6 & -0.0 & -1.0 & -0.6 & -8.0 & -0.1 & -0.9 & -0.9 \\
\bottomrule
\end{tabular}}
\end{table*}

\paragraph{\lrn{5}: The largest gains from extra capacity come from the level band, not from periodic structure.}
\emph{Evidence.} Replacing the linear backbone with an MLP lowers MSE by 13.5\% on Solar and 9.3\% on Weather (Table~\ref{tab:effects}); PatchTST improves on Cyc-Lin-CI by 20.4\% and 9.8\%. Most of this comes from the level band: MLP $\mlpLevelGainSolar$ and $\mlpLevelGainWeather$ points; PatchTST $\patchLevelGainSolar$ and $\patchLevelGainWeather$ points (Table~\ref{tab:bandgains}); on Solar PatchTST also improves the daily band ($\patchDailyGainSolar$). On the ETT datasets the MLP gives no consistent gain and is worse on ETTh1 and ETTm2. Across all cells the MLP is better in \cmpCycMLPCIvsCycLinCIwins{} (median $\cmpCycMLPCIvsCycLinCI\%$).
\emph{Caveat.} The gain does not track channel count: Traffic (862 channels) shows none. Our earlier statement that ``nonlinear capacity is vital for high-dimensional data'' is not supported.
\emph{Lesson.} On Solar and Weather the useful capacity predicts the \emph{mean level} of the next window (e.g.\ day versus night for Solar, slow weather regimes). Diagnosing this by band points to cheaper targeted fixes (a level-forecasting head, exogenous calendar features) that we have not yet tested.

\paragraph{\lrn{6}: Automatic channel mixing hurts more often than it helps.}
\emph{Evidence.} The automatic rule turned on channel mixing on seven of eight datasets. Forcing channel independence improves the MLP backbone in \cmpCycMLPCIvsCycMLPwins{} cells (median $\cmpCycMLPCIvsCycMLP\%$) and SOR in \cmpCycSORCIvsCycSORwins; for the linear and CNN backbones the difference is not consistent. The largest penalties are on Electricity and Traffic (+5\%), the datasets with the most channels.
\emph{Lesson.} With the current mixer and data sizes, channel independence is the safer default; the correlation threshold used to switch mixing on does not predict benefit.

\paragraph{\lrn{7}: Cheaper or fancier residual networks did not pay off.}
\emph{Evidence.} The 1D-CNN backbone uses \cnnParamsETT{} parameters on ETT versus \linParamsETT{} for the linear backbone, but is worse in all but \cmpCycCNNCIvsCycLinCIwins{} cells (median $\cmpCycCNNCIvsCycLinCI\%$) and has the worst average rank among trained models (\rankCycCNNCI). The spectrally conditioned patch mixer (SOR) is not better than the linear backbone (median $\cmpCycSORCIvsCycLinCI\%$, better in \cmpCycSORCIvsCycLinCIwins{}). Parameter counts are in Table~\ref{tab:params}.
\emph{Lesson.} Parameter reduction is only useful at matched accuracy; a 15\% error increase is a different operating point, not an efficiency gain.

\begin{table*}[t]
\centering
\caption{Zero-shot foundation models against trained models on identical non-overlapping test windows (seed 2024, MSE averaged over horizons). Bold/underline: best/second among the listed models. $^\dagger$Best trained model chosen per cell after the fact (optimistic reference). TTM-R1 was run only at $H{=}96$ and is discussed in the text.}
\label{tab:foundation}
\resizebox{\textwidth}{!}{\begin{tabular}{lccccccccc}
\toprule
& \multicolumn{5}{c}{Untrained / zero-shot} & \multicolumn{4}{c}{Trained on the target dataset} \\
\cmidrule(lr){2-6}\cmidrule(lr){7-10}
Dataset & S-Naive & Fourier top-3 & Chronos-T5 & Chronos-Bolt & TimesFM-2.5 & Linear & Cyc-Lin-CI & PatchTST & Best trained$^\dagger$ \\
\midrule
ETTh1 & 0.643 & 0.518 & 0.650 & 0.493 & 0.460 & 0.435 & \textbf{0.434} & \underline{0.434} & 0.429 \\
ETTh2 & 0.478 & 0.513 & 0.500 & 0.389 & 0.381 & \textbf{0.340} & \underline{0.351} & 0.353 & 0.332 \\
ETTm1 & 0.496 & 1.216 & 0.745 & 0.427 & 0.389 & 0.357 & \underline{0.352} & \textbf{0.351} & 0.347 \\
ETTm2 & 0.369 & 0.509 & 0.343 & 0.316 & 0.300 & 0.272 & \textbf{0.262} & \underline{0.267} & 0.260 \\
Electricity & 0.329 & 0.241 & 0.373 & \underline{0.163} & 0.168 & 0.167 & 0.167 & \textbf{0.162} & 0.157 \\
Solar & 0.351 & 1.695 & 0.842 & 0.738 & 0.343 & 0.274 & \underline{0.264} & \textbf{0.208} & 0.208 \\
Traffic & 1.226 & 0.683 & 1.138 & 0.476 & \textbf{0.397} & 0.436 & 0.423 & \underline{0.422} & 0.418 \\
Weather & 0.384 & 0.458 & 0.313 & 0.279 & 0.269 & 0.263 & \underline{0.256} & \textbf{0.233} & 0.231 \\
\bottomrule
\end{tabular}}
\end{table*}

\paragraph{\lrn{8}: Zero-shot foundation models beat seasonal naive but rarely a trained linear model.}
\emph{Evidence.} On the same windows, TimesFM-2.5 beats seasonal naive in \fmTimesFMTwoFiveVsSnWins{} cells but a trained Linear model in only \fmTimesFMTwoFiveVsLinWins{} (median $\fmTimesFMTwoFiveVsLinMed\%$). Chronos-Bolt beats Linear in \fmChronosBoltVsLinWins{} cells (median $\fmChronosBoltVsLinMed\%$); Chronos-T5 in none (median $\fmChronosTFiveVsLinMed\%$) and it is no better than seasonal naive in about half the cells; TTM-R1 beats Linear in \ttmVsLinWins{} cells at $H{=}96$. The exception is Traffic, where TimesFM is better than the best trained model at all \timesfmTrafficWins{} horizons (Table~\ref{tab:foundation}).
\emph{Caveat.} One seed and stride-$H$ windows (4--120 windows per cell); no fine-tuning.
\emph{Lesson.} On these benchmarks, a few minutes of training a linear map remains a strong baseline that zero-shot models need to be compared against, not only seasonal naive.

\begin{table}[t]
\centering
\caption{Effect of Fourier decoupling (top-3 rFFT base + model on the remainder) on zero-shot models: relative MSE change in \%, averaged over horizons (negative is better).}
\label{tab:fourier}
\resizebox{\columnwidth}{!}{\begin{tabular}{lcccc}
\toprule
Dataset & TimesFM-2.5 & Chronos-Bolt & Chronos-T5 & TTM-R1 ($H{=}96$) \\
\midrule
ETTh1 & +9 & +7 & -11 & +16 \\
ETTh2 & +25 & +27 & +14 & +55 \\
ETTm1 & +212 & +187 & +77 & +146 \\
ETTm2 & +68 & +63 & +51 & +109 \\
Electricity & +14 & +24 & -30 & -23 \\
Solar & +389 & +122 & +90 & +269 \\
Traffic & +22 & +12 & -20 & -2 \\
Weather & +65 & +51 & +35 & +83 \\
\midrule
cells improved & 0/32 & 1/32 & 13/32 & 2/8 \\
\bottomrule
\end{tabular}}
\end{table}

\paragraph{\lrn{9}: Fourier decoupling of pretrained models mostly hurts.}
\emph{Evidence.} Decoupling never improves TimesFM-2.5 (\fdecTimesFMTwoFiveImproved{} cells) and improves Chronos-Bolt in \fdecChronosBoltImproved{}. It helps Chronos-T5 in \fdecChronosTFiveImproved{} cells, mainly on Electricity, Traffic and ETTh1, and TTM in \fdecTTMImproved{} (Electricity, Traffic). The top-3 base on its own is worse than seasonal naive in \fourierBaseWorseThanSnaive{} cells, with very large errors on ETTm1 and Solar (Table~\ref{tab:fourier}).
\emph{Interpretation.} A three-component sinusoidal extrapolation is a poor forecast for most of these series; adding it only helps models whose own errors are larger still (the weakest model, Chronos-T5) on the most regular datasets. Our earlier ``inductive-bias rule'' (decoupling helps autoregressive models and hurts patch models) survives only in this narrow form, and the ``universal boundary'' we attributed to the Gibbs phenomenon on Weather is not specific to Weather: decoupling hurts on most datasets.
\emph{Caveat.} $K{=}3$ was fixed, not tuned. The size of the errors on ETTm1 and Solar should be checked against the implementation before drawing further conclusions (\todo{verify level handling of the Fourier base on 10/15-minute data}).

\begin{table}[t]
\centering
\caption{DSOF-SD spectral distillation with matched controls (test MSE, mean over three seeds). Control: same trainer and schedule with distillation weight zero. Warm: student after 5 epochs. Teacher: zero-shot TimesFM-2.5 on the full test set.}
\label{tab:dsof}
\resizebox{\columnwidth}{!}{\begin{tabular}{llccccccc}
\toprule
& & Teacher & \multicolumn{2}{c}{Linear student} & \multicolumn{4}{c}{SOR student} \\
\cmidrule(lr){4-5}\cmidrule(lr){6-9}
Data & $H$ & TimesFM & control & hard & warm & control & hard & soft \\
\midrule
ETTh1 & 96 & 0.395 & \textbf{0.373} & \underline{0.374} & 0.398 & 0.403 & 0.394 & 0.402 \\
 & 192 & 0.432 & \textbf{0.407} & \underline{0.408} & 0.470 & 0.499 & 0.457 & 0.467 \\
 & 336 & 0.451 & \textbf{0.432} & \underline{0.434} & 0.507 & 0.526 & 0.480 & 0.506 \\
 & 720 & \textbf{0.430} & 0.454 & \underline{0.453} & 0.492 & 0.543 & 0.506 & 0.516 \\
\midrule
ETTh2 & 96 & 0.314 & \underline{0.290} & \textbf{0.287} & 0.296 & 0.297 & 0.292 & 0.292 \\
 & 192 & 0.379 & 0.357 & 0.356 & \underline{0.345} & 0.345 & \textbf{0.343} & 0.348 \\
 & 336 & 0.404 & \textbf{0.376} & \underline{0.377} & 0.383 & 0.395 & 0.379 & 0.384 \\
 & 720 & 0.421 & 0.424 & \textbf{0.420} & 0.443 & 0.446 & \underline{0.420} & 0.422 \\
\bottomrule
\end{tabular}}
\end{table}

\paragraph{\lrn{10}: Spectral distillation acts as a regularizer for an overfitting student, not as a source of new information.}
\emph{Evidence.} For the linear student, distillation with hard gates matches the control (median $\dsofLINEARhardMed\%$, better in \dsofLINEARhardWins{}). For the SOR student it beats the control in \dsofSORhardWins{} cells (median $\dsofSORhardMed\%$), but the SOR control itself gets worse with the extra epochs (median $\sorControlVsWarmMed\%$ against the warm student), and the distilled SOR is still worse than the undistilled linear student (median $\sorHardVsLinCtrlMed\%$). The benefit-probe and consensus gates switched every band off in \benefitAllZeroPct\% and \consensusAllZeroPct\% of runs, which makes them identical to the control. The teacher is best only on ETTh1 at $H{=}720$ (Table~\ref{tab:dsof}).
\emph{Lesson.} Without a matched control, the earlier DSOF-SD results mixed the effect of distillation with the effect of 25 extra epochs. With the control, the measurable effect is early-stopping-like regularization.

\paragraph{\lrn{11}: The benchmark itself is reproducible; single-seed results are not.}
\emph{Evidence.} The identical configuration reproduces across two machines and GPU generations (validation MSE equal to within $10^{-9}$ in a spot check; \xmAgree{} cells within seed spread). The one large instability is Cyc-Lin on Traffic $H{=}720$ (\trafficSevenTwentyCycLin{} $\pm$ \trafficSevenTwentyCycLinStd{}, one diverging seed). The DLinear ``long-horizon breakdown'' on ETTh2 reported in our earlier draft does not exist under this protocol: DLinear reaches \dlinearEtthTwoSevenTwenty{} at $H{=}720$, among the best models (rFFT-Lin \rfftEtthTwoSevenTwenty, Cyc-Lin \cycLinEtthTwoSevenTwenty).
\emph{Lesson.} Most of the large effects we saw earlier were protocol effects; the true effects of the ingredients are of the order of 1--5\%.

\section{Claims Withdrawn from the Earlier Draft}
\label{sec:removed}
\begin{itemize}
    \item \textbf{``CycleNet-Pro removes DLinear's catastrophic breakdown on ETTh2 $H{=}720$ (0.4231 vs.\ 0.6050).''} Under one protocol DLinear scores \dlinearEtthTwoSevenTwenty{} and CycleNet-Pro (linear) \cycLinEtthTwoSevenTwenty{} (\lrn{11}).
    \item \textbf{``Mid-band ($f\in[0.1,0.4)$) FVU measures harmonic error; CycleNet cuts it by 31.8\%.''} With normalized-frequency bands the 24-step daily period, and the 96-, 144- and 168-step periods, all lie in the \emph{low} band; the mid band does not contain the cycles. Replaced by physical bands (Section~\ref{sec:diagnostics}).
    \item \textbf{``FVU$_\text{high}\ge 1$ shows systematic noise hallucination.''} An MSE-optimal forecast gives FVU $=1$ in an unpredictable band; this is expected behavior.
    \item \textbf{``CycleNet has the best spectral fidelity (lowest LSD).''} LSD can reward worse forecasts (Section~\ref{sec:diagnostics}); four different LSD definitions were also used across experiments.
    \item \textbf{``Nonlinear capacity is vital for high-dimensional data''; ``CNN gives 97.7\% fewer parameters with minimal accuracy loss.''} See \lrn{5} and \lrn{7}.
    \item \textbf{``Fourier decoupling: universal Gibbs boundary on Weather''; ``inductive-bias rule''}. See \lrn{9}.
    \item \textbf{Comparisons with published numbers at other lookbacks} (CycleNet/L at $L{=}96$, HDMixer). All comparisons are now run in the same harness.
\end{itemize}

\section{Threats to Validity}
\label{sec:threats}
\begin{itemize}
    \item \textbf{Statistics.} Wilcoxon tests over the 32 cells treat horizons of one dataset as independent, which they are not; we use the $p$-values descriptively. Per-window errors are stored for every run, so Diebold--Mariano tests per cell are possible and are the next step \todo{add DM tests for each claimed winner}.
    \item \textbf{Training budget.} PatchTST uses its published recipe (fixed learning rate, up to 100 epochs), all other models a 30-epoch budget with tuned learning rates. Part of the PatchTST advantage may come from this. The learning rate is chosen with one seed.
    \item \textbf{Scope.} Four of eight datasets are ETT variants; one lookback; no exogenous or calendar features; metrics in scaled space; foundation models without fine-tuning and on one seed.
    \item \textbf{Band definitions.} The 36\,h and 8\,h boundaries were fixed in advance; other choices shift shares between neighboring bands, though not the exactness of the split.
    \item \textbf{Open issues.} (a) Several \texttt{-CI} variants equal their non-CI twins on ETTh1, and Cyc-Fixed equals Cyc-Lin-CI wherever the automatic and fixed periods coincide; this follows from the configuration (the automatic rule disabled mixing on ETTh1) and is not a bug, but it means these cells are not independent evidence. (b) Traffic $H{=}720$ Cyc-Lin has one diverging seed. (c) TTM was only run at $H{=}96$. (d) The Fourier base error on ETTm1/Solar should be checked (\lrn{9}).
\end{itemize}

\section{Lessons and Next Steps}
\label{sec:lessons}
\begin{enumerate}
    \item \textbf{Measure the targetable error first.} The band split shows how much of a baseline's error sits in the band a new component addresses (\lrn{4}). On ETTh2, ETTm2 and Weather that ceiling is low for periodic models.
    \item \textbf{Change the function class or the target, not just the basis.} Linear transforms before a linear head are equivalent (\lrn{2}).
    \item \textbf{Treat the period as a hyperparameter selected on validation error} (\lrn{3}).
    \item \textbf{Default to channel independence} with the current mixers (\lrn{6}).
    \item \textbf{Compare against a trained linear model and seasonal naive}, in the same harness, with several seeds (\lrn{1}, \lrn{8}, \lrn{11}).
    \item \textbf{Promising directions supported by the data:} a level/regime head for Solar and Weather (\lrn{5}); validated period selection including weekly cycles (\lrn{3}); per-window Diebold--Mariano tests; equal training budgets for all models.
\end{enumerate}

\section{Conclusion}
Under a single controlled protocol, frequency-domain modeling helps lightweight forecasters in one specific situation: when the series has a stable, correctly identified cycle that carries a large share of the error, as on Solar, Traffic, ETTm1 and (with the weekly period) Electricity. It does not help as a change of basis alone, it cannot help where the error is in the level and slow bands, and its benefits are of the same order as the effects of period choice, channel mixing and training budget. Pretrained models with Fourier decoupling and spectral distillation did not change this picture. We hope the corrected diagnostics, the matched-control design and the withdrawn claims are as useful to others as the positive findings.

\bibliographystyle{plain}
\bibliography{references}

\clearpage
\appendix
\section{Full Results}
\begin{table*}[h]
\centering
\caption{Test MSE (mean over seeds 2024--2026) for all cells. Bold/underline: best/second among trained models in the row. S-Naive is the untrained reference.}
\label{tab:full}
\resizebox{\textwidth}{!}{\begin{tabular}{llccccccccccccc}
\toprule
Data & $H$ & S-Naive & Linear & DLinear & rFFT-Lin & DCT-Lin & Haar-Lin & STFT-Lin & Cyc-Fixed & Cyc-Lin-CI & Cyc-MLP-CI & Cyc-SOR-CI & Cyc-CNN-CI & PatchTST \\
\midrule
ETTh1 & 96 & 0.512 & 0.376 & 0.376 & \underline{0.372} & 0.373 & 0.374 & 0.393 & 0.372 & 0.372 & 0.381 & 0.391 & 0.431 & \textbf{0.371} \\
 & 192 & 0.581 & 0.414 & 0.410 & 0.417 & 0.407 & 0.411 & 0.423 & \textbf{0.407} & \underline{0.407} & 0.416 & 0.430 & 0.439 & 0.410 \\
 & 336 & 0.650 & 0.435 & 0.434 & 0.439 & 0.432 & 0.435 & 0.442 & \textbf{0.430} & \underline{0.430} & 0.436 & 0.446 & 0.446 & 0.438 \\
 & 720 & 0.655 & 0.445 & 0.478 & \textbf{0.442} & \underline{0.444} & 0.445 & 0.447 & 0.459 & 0.459 & 0.472 & 0.501 & 0.464 & 0.457 \\
\midrule
ETTh2 & 96 & 0.391 & 0.277 & 0.283 & \underline{0.274} & \textbf{0.273} & 0.276 & 0.303 & 0.279 & 0.279 & 0.298 & 0.290 & 0.362 & 0.277 \\
 & 192 & 0.482 & 0.343 & 0.341 & \underline{0.338} & \textbf{0.337} & 0.340 & 0.347 & 0.343 & 0.343 & 0.352 & 0.357 & 0.390 & 0.350 \\
 & 336 & 0.532 & 0.373 & 0.372 & 0.375 & \textbf{0.363} & 0.370 & \underline{0.369} & 0.395 & 0.395 & 0.376 & 0.379 & 0.387 & 0.380 \\
 & 720 & 0.525 & 0.395 & \underline{0.393} & \textbf{0.391} & 0.400 & 0.393 & 0.397 & 0.419 & 0.419 & 0.408 & 0.415 & 0.438 & 0.396 \\
\midrule
ETTm1 & 96 & 0.423 & 0.306 & 0.302 & 0.302 & 0.302 & 0.318 & 0.320 & 0.301 & 0.301 & \textbf{0.286} & 0.295 & 0.361 & \underline{0.287} \\
 & 192 & 0.463 & 0.344 & 0.345 & 0.337 & 0.336 & 0.341 & 0.344 & 0.335 & 0.335 & \textbf{0.328} & 0.332 & 0.380 & \underline{0.331} \\
 & 336 & 0.496 & 0.377 & 0.376 & 0.373 & 0.374 & 0.377 & 0.382 & 0.368 & 0.368 & \textbf{0.363} & \underline{0.365} & 0.401 & 0.368 \\
 & 720 & 0.574 & 0.432 & 0.429 & 0.429 & 0.429 & 0.431 & 0.442 & \underline{0.425} & 0.425 & 0.429 & 0.445 & 0.438 & \textbf{0.418} \\
\midrule
ETTm2 & 96 & 0.263 & 0.166 & 0.165 & 0.165 & 0.164 & 0.165 & 0.172 & \textbf{0.162} & \underline{0.162} & 0.170 & 0.173 & 0.246 & 0.166 \\
 & 192 & 0.321 & 0.222 & 0.221 & 0.221 & 0.220 & 0.221 & 0.238 & \underline{0.216} & \textbf{0.216} & 0.229 & 0.230 & 0.276 & 0.224 \\
 & 336 & 0.376 & 0.295 & 0.297 & 0.274 & 0.273 & 0.274 & 0.292 & \textbf{0.271} & \underline{0.271} & 0.281 & 0.283 & 0.312 & 0.276 \\
 & 720 & 0.471 & 0.371 & 0.368 & 0.368 & 0.367 & 0.369 & 0.370 & \underline{0.365} & 0.366 & 0.366 & 0.369 & 0.393 & \textbf{0.363} \\
\midrule
Electricity & 96 & 0.321 & 0.140 & 0.141 & 0.140 & 0.140 & 0.140 & 0.151 & \textbf{0.129} & 0.140 & 0.131 & 0.135 & 0.234 & \underline{0.131} \\
 & 192 & 0.304 & 0.154 & 0.154 & 0.154 & 0.154 & 0.154 & 0.166 & \textbf{0.144} & 0.154 & \underline{0.148} & 0.155 & 0.244 & 0.150 \\
 & 336 & 0.327 & 0.172 & 0.172 & 0.171 & 0.171 & 0.171 & 0.175 & \textbf{0.161} & 0.171 & \underline{0.166} & 0.177 & 0.197 & 0.167 \\
 & 720 & 0.367 & 0.210 & 0.210 & 0.210 & 0.211 & 0.210 & 0.212 & \textbf{0.199} & 0.210 & \underline{0.203} & 0.215 & 0.285 & 0.203 \\
\midrule
Solar & 96 & 0.290 & 0.226 & 0.226 & 0.226 & 0.226 & 0.226 & 0.240 & 0.220 & 0.220 & \underline{0.197} & 0.197 & 0.262 & \textbf{0.177} \\
 & 192 & 0.327 & 0.257 & 0.257 & 0.257 & 0.257 & 0.259 & 0.268 & 0.247 & 0.247 & 0.217 & \underline{0.212} & 0.281 & \textbf{0.194} \\
 & 336 & 0.372 & 0.280 & 0.280 & 0.279 & 0.279 & 0.280 & 0.288 & 0.271 & 0.271 & 0.223 & \underline{0.218} & 0.287 & \textbf{0.211} \\
 & 720 & 0.376 & 0.280 & 0.279 & 0.279 & 0.279 & 0.279 & 0.284 & 0.268 & 0.268 & 0.231 & \underline{0.219} & 0.281 & \textbf{0.217} \\
\midrule
Traffic & 96 & 1.217 & 0.410 & 0.410 & 0.410 & 0.410 & 0.410 & 0.418 & 0.400 & 0.400 & \underline{0.393} & 0.402 & 0.622 & \textbf{0.384} \\
 & 192 & 1.093 & 0.423 & 0.423 & 0.423 & 0.423 & 0.423 & 0.433 & 0.415 & 0.414 & \textbf{0.408} & 0.415 & 0.629 & \underline{0.409} \\
 & 336 & 1.150 & 0.435 & 0.435 & 0.435 & 0.435 & 0.435 & 0.442 & \underline{0.426} & 0.426 & \textbf{0.425} & 0.430 & 0.532 & 0.429 \\
 & 720 & 1.185 & 0.463 & 0.463 & 0.463 & 0.464 & 0.463 & 0.469 & \underline{0.453} & \textbf{0.447} & 0.462 & 0.466 & 0.647 & 0.462 \\
\midrule
Weather & 96 & 0.317 & 0.175 & 0.175 & 0.174 & 0.175 & 0.175 & 0.199 & 0.178 & 0.178 & \underline{0.148} & 0.148 & 0.234 & \textbf{0.146} \\
 & 192 & 0.343 & 0.218 & 0.217 & 0.218 & 0.218 & 0.218 & 0.228 & 0.212 & 0.212 & \textbf{0.190} & 0.193 & 0.262 & \underline{0.190} \\
 & 336 & 0.383 & 0.271 & 0.265 & 0.265 & 0.267 & 0.266 & 0.274 & 0.260 & 0.260 & \underline{0.242} & 0.250 & 0.295 & \textbf{0.242} \\
 & 720 & 0.443 & 0.333 & 0.333 & 0.333 & 0.333 & 0.336 & 0.337 & 0.327 & 0.327 & \underline{0.319} & 0.324 & 0.348 & \textbf{0.316} \\
\bottomrule
\end{tabular}}
\end{table*}

\begin{table*}[h]
\centering
\caption{Trainable parameters per dataset.}
\label{tab:params}
\resizebox{\textwidth}{!}{\begin{tabular}{lrrrrrrrr}
\toprule
Model & ETTh1 & ETTh2 & ETTm1 & ETTm2 & Electricity & Solar & Traffic & Weather \\
\midrule
Linear & 32,352 & 32,352 & 32,352 & 32,352 & 32,352 & 32,352 & 32,352 & 32,352 \\
DLinear & 64,608 & 64,608 & 64,608 & 64,608 & 64,608 & 64,608 & 64,608 & 64,608 \\
STFT-Lin & 57,120 & 57,120 & 57,120 & 57,120 & 57,120 & 57,120 & 57,120 & 57,120 \\
Cyc-Fixed & 32,534 & 32,534 & 33,038 & 33,038 & 86,922 & 52,354 & 178,892 & 35,418 \\
Cyc-Lin-CI & 32,534 & 32,534 & 33,038 & 33,318 & 40,698 & 52,354 & 199,580 & 35,418 \\
Cyc-Lin & 32,534 & 32,601 & 33,105 & 33,385 & 143,899 & 71,191 & 943,917 & 35,869 \\
Cyc-MLP-CI & 221,974 & 221,974 & 222,478 & 222,758 & 230,138 & 241,794 & 389,020 & 224,858 \\
Cyc-SOR-CI & 76,668 & 76,668 & 77,172 & 77,452 & 84,832 & 96,488 & 243,714 & 79,552 \\
Cyc-CNN-CI & 5,559 & 5,559 & 6,063 & 6,343 & 13,723 & 25,379 & 172,605 & 8,443 \\
Cyc-CNN & 5,559 & 5,626 & 6,130 & 6,410 & 116,924 & 44,216 & 916,942 & 8,894 \\
PatchTST & 81,742 & 81,742 & 921,198 & 921,198 & 921,826 & 921,458 & 922,908 & 921,226 \\
\bottomrule
\end{tabular}}
\end{table*}

\end{document}
